% CV de Baye Lahad Mbacke -- compiler avec : pdflatex Resume_PHD.tex
\documentclass[11pt,a4paper]{article}
\usepackage[utf8]{inputenc}
\usepackage[T1]{fontenc}
\usepackage{lmodern}
\usepackage[margin=1.8cm,top=1.5cm,bottom=1.5cm]{geometry}
\usepackage{xcolor}
\usepackage{tabularx}
\usepackage[hidelinks,pdftitle={Baye Lahad Mbacke -- CV},pdfauthor={Baye Lahad Mbacke}]{hyperref}

\definecolor{accent}{HTML}{0066CC}
\pagestyle{empty}
\setlength{\parindent}{0pt}

\newcommand{\cvsection}[1]{%
  \vspace{10pt}{\Large\color{accent}#1}\par\vspace{-6pt}%
  {\color{accent}\rule{\linewidth}{0.8pt}}\par\vspace{4pt}}

% \cventry{institution}{lieu}{intitulé}{dates}
\newcommand{\cventry}[4]{%
  \begin{tabularx}{\linewidth}{@{}X r@{}}
    \textbf{#1} & #2 \\
    \textit{#3} & \textit{#4}
  \end{tabularx}\par}

\newenvironment{cvitems}{%
  \begin{list}{--}{\setlength{\leftmargin}{1.5em}\setlength{\itemsep}{1pt}
  \setlength{\parsep}{0pt}\setlength{\topsep}{2pt}}}{\end{list}\vspace{4pt}}

\begin{document}

\begin{center}
  {\fontsize{24}{28}\selectfont\color{accent} Baye Lahad Mbacke}\\[6pt]
  Doctorant en Informatique -- Nantes Université, LS2N\\[2pt]
  8 boulevard Mendès France, Nantes, France \textbullet{} 06 01 32 80 22\\[2pt]
  \href{mailto:lahadmbacke00@gmail.com}{lahadmbacke00@gmail.com} \textbullet{}
  \href{https://lahadmbacke.github.io/Lahad_Mbacke/}{Portfolio} \textbullet{}
  \href{https://linkedin.com/in/baye-lahad-mbacke}{LinkedIn} \textbullet{}
  \href{https://github.com/LahadMbacke}{GitHub} \textbullet{}
  \href{https://huggingface.co/Lahad}{Hugging Face}
\end{center}

\vspace{4pt}
Doctorant à Nantes Université depuis mars 2026, au sein du Laboratoire des Sciences du Numérique
de Nantes (LS2N). Mes travaux portent sur des méthodes efficaces et frugales d'extraction et de
structuration automatiques d'informations sur les événements culturels, afin de connecter les
acteurs des industries culturelles et créatives à leurs publics.

\cvsection{Formation Académique}

\cventry{Nantes Université -- LS2N (équipe TALN) \& LIA}{Nantes, France}
        {Doctorat en Informatique (École doctorale MaSTIC)}{Mars 2026--Aujourd'hui}
\textit{Extraction et structuration automatique d'événements culturels : une approche efficace et
frugale pour connecter les acteurs des ICC et leurs publics}
\begin{cvitems}
  \item Direction : Richard Dufour (LS2N, Nantes Université) et Vincent Labatut (LIA, Avignon Université)
  \item Partenaire industriel : ideactiv (plateforme française d'indexation culturelle)
  \item Détection des listes d'événements sur les sites d'acteurs culturels, identification des
        événements, extraction et consolidation de leurs informations (titre, date, lieu, prix,
        réservation, etc.), sans hallucination ni recours à des services tiers coûteux
\end{cvitems}

\cventry{Nantes Université}{Nantes, France}
        {Master ATAL (Apprentissage et Traitement Automatique de la Langue)}{2023--2025}
Spécialisation en techniques avancées de TAL et d'apprentissage automatique\par\vspace{6pt}

\cventry{Université de Strasbourg}{Strasbourg, France}{L3 en Informatique}{2022--2023}
Approfondissement en algorithmique et programmation\par\vspace{6pt}

\cventry{Université Cheikh Anta Diop de Dakar}{Dakar, Sénégal}
        {Licence en Mathématiques-Informatique}{2018--2022}
Formation aux fondamentaux en sciences mathématiques et informatiques

\cvsection{Expérience Professionnelle}

\cventry{Ouest-France}{Rennes, France}
        {Assistant Data Scientist -- Recherche \& Développement (alternance)}{Sept. 2024--Sept. 2025}
\begin{cvitems}
  \item Geotagging : identification de références de localisation à partir de bribes d'adresse
        dans des articles journalistiques, et rattachement à des tags de lieux
  \item Extraction de connaissances : reconnaissance et mise en relation d'entités nommées
  \item Désambiguïsation d'entités géographiques (adresse, place, monument, ville, etc.)
\end{cvitems}

\cventry{Laboratoire des Sciences du Numérique de Nantes (LS2N)}{Nantes, France}
        {Stagiaire Recherche}{Janv. 2024--Mai 2024}
\begin{cvitems}
  \item Comparaison de systèmes de reconnaissance d'entités nommées sur des articles de presse
  \item Apprentissage automatique pour l'extraction d'informations clés
  \item Annotation et analyse de données structurées et non structurées
\end{cvitems}

\cvsection{Projets Académiques}

\cventry{Nantes Université}{Nantes, France}
        {Compétition DEFT2013 -- Classification de recettes de cuisine}{Févr. 2024--Mars 2024}
\begin{cvitems}
  \item Modèles de classification de recettes (plat principal, entrée, dessert) à partir des
        titres et des instructions, par des méthodes de TAL et d'apprentissage automatique
\end{cvitems}

\cventry{Université de Strasbourg}{Strasbourg, France}
        {Auto-encodeurs et classification}{Avr. 2023--Mai 2023}
\begin{cvitems}
  \item Évaluation de l'intérêt des auto-encodeurs pour la classification de données tabulaires
  \item Analyse comparative des performances avec et sans compression dans un espace latent réduit
\end{cvitems}

\cvsection{Compétences Techniques}
\begin{tabularx}{\linewidth}{@{}l X@{}}
  \textbf{Machine \& Deep Learning} & PyTorch, Scikit-learn, LLM, Fine-tuning, Quantization, Distillation \\
  \textbf{TAL} & spaCy, NLTK, LangChain, RAG, Agents, Reconnaissance d'entités nommées \\
  \textbf{Programmation} & Python, R, Java, C/C++ \\
  \textbf{Bases de données} & SQL, MySQL, Oracle, PostgreSQL, MongoDB
\end{tabularx}

\cvsection{Langues}
\textbf{Wolof} : langue maternelle \quad\textbullet\quad
\textbf{Français} : courant \quad\textbullet\quad
\textbf{Anglais} : B2

\cvsection{Centres d'Intérêt}
\textbf{Lecture} : journaux, actualités technologiques, veille informatique \quad\textbullet\quad
\textbf{Sport} : football

\end{document}
